\section{绪论}
\label{sec:introduction}

\subsection{研究背景与意义}

多功能雷达/电子对抗网络是现代电子战的核心组成部分。随着电磁环境日趋复杂，单一雷达已无法满足广域覆盖、高可靠探测和有效电子压制的综合需求，多雷达协同组网成为必然趋势。有效的雷达空间部署直接影响网络的两项核心效能：探测感知能力（覆盖率）和电子干扰压制能力（最小等效干扰强度）。

近年来，无人机（UAV）平台的广泛应用为空地协同电子对抗网络带来了新的机遇。空中节点通过空-地（A2G）链路传播，具有近似自由空间的传播特性（路径损耗指数$\alpha \approx 2.0$--2.5），可实现广域覆盖；地面节点通过地-地（G2G）链路传播，受地形遮挡和多径效应影响显著（$\alpha \approx 3.5$--4.0），但可提供持续稳定的干扰压制能力。如何平衡这两类节点的部署以最大化覆盖效能，同时保证干扰压制能力，是空地协同电子对抗网络部署优化的核心问题。

然而，传统的雷达部署方法面临两个关键挑战：

\textbf{边界效应问题。}在复杂区域部署雷达网络时，区域边缘的探测概率显著低于中心区域。Han等~\cite{Han2025}指出，传统方法在区域边界处的覆盖率通常比中心区域低15\%--20\%。这一现象的根本原因在于优化搜索空间与物理部署空间之间的不均匀映射，导致边界区域被欠采样。

\textbf{传播模型简化问题。}现有方法通常忽略空地链路的传播差异，采用统一的路径损耗模型。在空地协同场景中，空中节点和地面节点的传播环境截然不同：A2G链路接近自由空间，传播距离远且衰减慢；G2G链路受地面杂波、多径和遮挡影响，衰减显著增大。忽略这一差异将导致覆盖计算的系统偏差。

针对上述挑战，本文基于Han等~\cite{Han2025}提出的MOPSO-DT算法框架，引入空地异构传播模型并优化算法参数，系统地研究空地协同电子对抗网络在复杂区域的多目标部署优化问题。

\subsection{国内外研究现状}

\textbf{雷达网络部署优化。}雷达网络部署优化是电子对抗和防空预警领域的经典问题。早期工作以单目标优化为主，采用遗传算法（GA）或模拟退火等元启发式方法求解最大覆盖率部署方案。Kirubarajan等率先将进化算法系统地应用于该领域。近年来，多目标优化方法被广泛引入：Chen等~\cite{Chen2023}使用NSGA-II同时优化雷达网络覆盖率和部署成本；Kumar和Sharma~\cite{Kumar2022}研究了区域分解方法在雷达部署中的应用，通过将复杂区域分解为子区域来降低搜索难度。

然而，现有方法普遍忽视了两个关键问题：一是区域边界处的探测概率系统性下降（即"边界效应"），根源在于搜索空间与物理部署空间之间的不均匀映射导致边界区域被欠采样——Han等~\cite{Han2025}的实验表明传统方法在区域边界处的覆盖率通常比中心区域低15\%--20\%；二是空地链路传播特性差异对覆盖计算的影响——绝大多数方法采用统一的路径损耗模型，忽略了空中节点（A2G链路，路径损耗指数$\alpha \approx 2.0$--2.5）和地面节点（G2G链路，$\alpha \approx 3.5$--4.0）在传播环境上的本质区别。

\textbf{多目标进化算法。}进化多目标优化（EMO）是计算智能领域最活跃的研究分支之一。Deb等~\cite{Deb2002}提出的NSGA-II通过快速非支配排序和拥挤距离比较算子，以$O(M N^2)$的时间复杂度高效求解多目标优化问题，成为该领域引用最广泛的基准算法。Zitzler等~\cite{Zitzler2001}提出的SPEA2采用强度Pareto适应度分配和$k$-近邻密度估计，在多样性维护上表现出色。Zhang和Li~\cite{Zhang2007}提出的MOEA/D将多目标问题分解为一组标量子问题协同优化，在高维目标场景下具有独特优势。Hua等~\cite{Hua2021}对非规则Pareto前沿的处理方法进行了系统综述，Vachhani等~\cite{Vachhani2015}对MOEA的主流方法进行了分类和比较。

在雷达部署这一具体应用领域，Yan等~\cite{Yan2024}利用改进NSGA-II求解组网雷达搜索-跟踪联合部署问题，验证了Pareto优化框架的有效性。Han等~\cite{Han2025}提出的MOPSO-DT算法在MOPSO基础上引入了区域凸分解和坐标变换步骤以抑制边界效应，但未考虑空地异构传播特性。

\textbf{MOPSO及领导选择策略。}多目标粒子群优化（MOPSO）由Coello和Lechuga于2002年首次提出，Coello等~\cite{Coello2004}随后给出了MOPSO的规范化框架。Raquel和Naval~\cite{Raquel2005}首次将拥挤距离引入MOPSO（MOPSO-CD），用于外部档案维护和全局最优选择，开创了"拥挤距离 + 外部归档"的设计范式。在全局最优选择方面，Mostaghim和Teich~\cite{Mostaghim2003}提出了基于sigma值的几何引导策略。2023--2025年间，该方向取得了显著进展：Liu等~\cite{Liu2023}的TMMOPSO采用超锥域自适应选择；Wang和Zhou~\cite{Wang2024}的PPSO利用自组织映射预测领导者；Zhang等~\cite{Zhang2024}的MOPSO/CP基于局部理想点构造收敛贡献评估器。上述进展表明，\textbf{全局最优选择策略是MOPSO性能的关键决定因素}，合理的选择策略可同时提升收敛速度和Pareto前沿的分布质量。

\textbf{空地传播模型。}国际电信联盟（ITU-R）在建议书~\cite{ITUP528}中给出了航空移动业务的传播曲线，建议A2G链路路径损耗指数约为2.0--2.5。建议书~\cite{ITUP452}给出了地面站间干扰评估方法，G2G链路在典型城市环境下路径损耗指数约为3.5--4.0。Rappaport~\cite{Rappaport2002}系统总结了不同传播环境下的路径损耗模型。本文采用区分A2G和G2G链路的异构传播模型，对空中节点和地面节点分别设定不同的路径损耗指数，更准确地刻画空地协同网络的真实传播特性。

\textbf{计算几何与区域分解。}Hertel和Mehlhorn~\cite{Hertel1983}提出的快速三角剖分算法可在$O(n \log n)$时间内完成简单多边形的三角剖分。de Berg等~\cite{DeBerg2008}系统总结了计算几何中的凸分解和坐标变换方法。Chazelle~\cite{Chazelle1982}的多边形切割定理为复杂区域的空洞消除提供了理论基础。本文沿袭MOPSO-DT的Hertel-Mehlhorn区域凸分解和垂直交点坐标变换方法，将归一化搜索空间中的候选解一一映射到非凸多边形区域的合法物理位置。

综合以上分析，现有研究缺乏同时处理（1）复杂区域边界效应抑制、（2）空地异构传播差异建模、（3）覆盖率与干扰压制效能联合优化、（4）基于拥挤距离的全局最优选择策略的统一框架。本文旨在填补这一空白，并通过系统的参数调优和消融实验验证各组件的独立贡献。

\subsection{研究内容与创新点}

本文围绕以下三个核心问题展开研究：

\begin{itemize}
    \item \textbf{Q1（边界效应消除）：}如何通过区域分解和坐标变换消除雷达网络部署中的边界效应？
    \item \textbf{Q2（异构传播建模）：}如何建立准确反映空地传播差异的优化模型？
    \item \textbf{Q3（多目标权衡）：}如何在覆盖率（ECR）和干扰压制效能（$J_{\text{min}}$）之间取得最优权衡？
\end{itemize}

\textbf{本文主要创新点：}

\begin{enumerate}
    \item 实现Hertel-Mehlhorn区域凸分解与垂直交点坐标变换，有效抑制部署优化中的边界效应
    \item 引入空地异构传播模型（A2G: $\alpha=2.0$, G2G: $\alpha=4.0$），更准确地刻画空地协同网络的真实传播特性
    \item 对MOPSO-DT进行系统参数调优：standard惯性权重策略结合crowding拥挤度全局最优选择，在超体积（HV）上提升5.9\%，Pareto解空间分布均匀性（Spacing）改善54\%
    \item 在两类实验场景中定量观察到ECR与最小等效干扰强度$J_{\text{min}}$之间的强负相关关系（$r < -0.93$），为有限资源条件下覆盖率与压制效能的权衡决策提供实验证据
\end{enumerate}

\subsection{论文结构安排}

本文共分为五章。第一章为绪论，介绍研究背景、相关工作和本文贡献。第二章介绍相关理论基础，包括雷达网络模型、多目标优化理论、粒子群优化算法和计算几何基础。第三章详细描述问题建模与MOPSO-DT算法设计。第四章给出实验设置、结果分析和讨论。第五章总结全文并提出未来研究方向。
